\documentclass[12pt]{article}
\usepackage[T1]{fontenc}
\usepackage[utf8]{inputenc}
\usepackage{lmodern}
\usepackage{textcomp}
\usepackage{enumitem}
\usepackage{geometry}
\usepackage{graphicx}
\usepackage{amsmath, amssymb}
\geometry{margin=1in}
\begin{document}
\begin{center}
Physics - Undergraduate Level Assessment 14 \
Total Marks: 40 \
Answer all questions.
\end{center}

\section*{Multiple Choice (10 marks)}
\begin{enumerate}[label=\arabic*.]
\item A rod measures 1.00 m in its rest system. How fast must an observer move parallel to the rod to measure its length to be 0.80 m?
\begin{enumerate}[label=(\alph*)]
    \item 0.50 c
    \item 0.60 c
    \item 0.70 c
    \item 0.80 c
\end{enumerate}
\item For which of the following thermodynamic processes is the increase in the internal energy of an ideal gas equal to the heat added to the gas?
\begin{enumerate}[label=(\alph*)]
    \item Constant temperature
    \item Constant volume
    \item Constant pressure
    \item Adiabatic
\end{enumerate}
\item In the diamond structure of elemental carbon, the nearest neighbors of each C atom lie at the corners of a
\begin{enumerate}[label=(\alph*)]
    \item square
    \item hexagon
    \item cube
    \item tetrahedron
\end{enumerate}
\item The rest mass of a particle with total energy 5.0 GeV and momentum 4.9 GeV/c is approximately
\begin{enumerate}[label=(\alph*)]
    \item 0.1 GeV/$c^2$
    \item 0.2 GeV/$c^2$
    \item 0.5 GeV/$c^2$
    \item 1.0 GeV/$c^2$
\end{enumerate}
\item The energy required to remove both electrons from the helium atom in its ground state is 79.0 eV. How much energy is required to ionize helium (i.e., to remove one electron)?
\begin{enumerate}[label=(\alph*)]
    \item 24.6 eV
    \item 39.5 eV
    \item 51.8 eV
    \item 54.4 eV
\end{enumerate}
\end{enumerate}

\section*{True / False (10 marks)}
\textit{Answer True or False}
\begin{enumerate}[label=\arabic*.]
\setcounter{enumi}{5}
\item True or False: When the voltage across a resistor is doubled, the rate of energy dissipation increases to 4 W due to the power formula P = $V^2$/R.
\item True or False: The electric displacement current through a surface is proportional to the time integral of the magnetic flux through that surface.
\item True or False: Fermions have antisymmetric wave functions and obey the Pauli exclusion principle, preventing them from occupying the same quantum state.
\item True or False: The speed of light in a nonmagnetic dielectric material with a dielectric constant of 4.0 is 3.0 * $10^8$ m/s.
\item True or False: A dye laser is the most effective type of laser for spectroscopy across a broad range of visible wavelengths.
\end{enumerate}

\section*{Long Form Response (20 marks)}
\textit{Answer each question with a clear and complete explanation.}
\begin{enumerate}[label=\arabic*.]
\setcounter{enumi}{10}
\item Discuss the behavior of unpolarized light passing through two ideal linear polarizers with their transmission axes at a 45-degree angle. Calculate the transmitted light intensity as a percentage of the incident intensity and explain the underlying principles.
\item Discuss the role of various ions as dopants in germanium semiconductors, specifically identifying which ions are ineffective in creating n-type semiconductors and why.
\end{enumerate}

\vfill
\noindent\textit{End of Paper}
\end{document}